\documentclass[12pt,a4paper]{article}

\usepackage[margin=2.5cm]{geometry}
\usepackage{amsmath,amssymb}
\usepackage{graphicx}
\usepackage{booktabs}
\usepackage{natbib}
\usepackage{hyperref}
\usepackage{setspace}
\usepackage{microtype}
\usepackage{float}
\usepackage{caption}
\usepackage{subcaption}
\usepackage{array}
\usepackage{multirow}
\usepackage{xcolor}
\usepackage{titlesec}
\usepackage{abstract}

\doublespacing
\bibpunct{(}{)}{;}{a}{,}{,}

\hypersetup{
    colorlinks=true,
    linkcolor=blue!70!black,
    citecolor=blue!70!black,
    urlcolor=blue!70!black
}

\title{\textbf{Urbanism Quality, Housing Prices, and Effective Housing Supply:\\
Evidence from Israeli Cities}\thanks{%
  Data: Israeli Tax Authority real-estate transaction records (2018--2023);
  OpenStreetMap road network and amenity data retrieved via Overture Maps (2024);
  Israeli Central Bureau of Statistics statistical areas (2022 edition).
  All code and replication materials are available at
  \texttt{github.com/mulysan/housing}.}}

\author{%
  Mulya Sanderovitch\\
  \small\textit{Working Paper, March 2026}
}

\date{}

\begin{document}

\maketitle
\thispagestyle{empty}

%----------------------------------------------------------------------
\begin{abstract}
\noindent
This paper links observable urbanism metrics—street-network density, junction
density, amenity density, and composite walkability—to residential apartment
prices across 21 Israeli cities, using administrative transaction data for
2018–2023 and street-network/amenity data from OpenStreetMap.  We find that
street density (km/km$^2$) and amenity density (POIs/km$^2$) are the strongest
cross-city predictors of median apartment prices, with Pearson correlations of
0.63 and 0.64 respectively (both $p < 0.01$), while a composite walkability
index carries a Pearson correlation of 0.44 ($p < 0.05$).  Dead-end ratio and
circuity are not significantly associated with prices at the city level.  We
then embed these findings in a broader theoretical framework, drawing on the
hedonic price literature and on spatial equilibrium models of housing supply.
We argue that the observed capitalization of neighborhood walkability and
amenity access is precisely the micro-level evidence needed to construct a
\textit{quality-adjusted} or ``effective'' housing supply measure—analogous to
efficiency-weighted labor in labor economics—in which each dwelling unit
contributes to the aggregate stock in proportion to its hedonic-predicted value.
Accounting for this quality margin implies that the Israeli housing shortage is
more severe than raw unit counts suggest, because a disproportionate share of
the existing stock is located in low-connectivity, low-amenity urban
environments that provide fewer ``housing services'' per physical unit.

\bigskip\noindent
\textbf{JEL codes:} R21, R31, R41, R52, D40\\
\textbf{Keywords:} hedonic prices, walkability, street-network connectivity,
amenity density, effective housing supply, spatial equilibrium, Israel
\end{abstract}

\newpage
\tableofcontents
\newpage

%======================================================================
\section{Introduction}
\label{sec:intro}
%======================================================================

Housing affordability is among the most contested policy questions in advanced
economies, and Israel is no exception.  Between 2008 and 2023, real apartment
prices in the Tel Aviv metropolitan area roughly doubled, even after controlling
for income growth \citep{boi2023housing}.  Public debate has overwhelmingly
focused on the \textit{quantity} dimension of supply: how many new units are
built per year, how quickly permits are issued, and how zoning reform can
accelerate construction.  Largely absent from this debate is the
\textit{quality} dimension: dwellings are not homogeneous, and a unit built in
a walkable, transit-accessible neighborhood with dense retail and services
provides fundamentally different ``housing services'' than a physically
identical unit in a car-dependent, low-amenity suburb.

This paper makes two contributions.  First, it provides the first systematic
empirical analysis linking quantitative urbanism metrics—derived from
OpenStreetMap road-network graphs and point-of-interest data—to residential
apartment prices across a broad cross-section of Israeli cities.  Our data
cover 21 cities, approximately 200,000 apartment transactions over 2018–2023,
and six urbanism metrics computed at the city level from Overture Maps data.
We find robust positive associations between street-network density, amenity
density, and median apartment prices.  These associations survive both Pearson
and Spearman specifications and are economically large: a city at the 75th
percentile of amenity density commands a median price roughly 60\% higher than
a city at the 25th percentile.

Second, we connect these micro-level correlations to the macro-level debate on
housing supply.  Building on the concept of ``effective labor supply'' in labor
economics—where heterogeneous workers are aggregated into efficiency units using
relative wages as weights \citep{katz1992changes}—we argue for an analogous
concept of \textit{effective housing supply}.  Just as a high-skill worker
earning twice the median wage contributes two efficiency units of labor, a
dwelling in a walkable, well-connected neighborhood whose hedonic-predicted
value is twice the market average contributes two ``housing service units'' to
the aggregate stock.  The lit\-erature reviewed in Section~\ref{sec:lit}
provides the micro-foundations: the hedonic capitalization of walkability,
street connectivity, transit access, and amenity density quantifies exactly how
much more ``service'' a well-located unit provides.

The distinction between quantity and quality supply matters for at least three
reasons.  First, it changes how we measure the housing deficit.  If Israel's
housing shortage is partly a \textit{quality} shortage—too few units in
walkable, amenity-rich areas—then policies that simply add units at the urban
periphery may fail to close the effective gap.  Second, it changes how we
interpret supply elasticities.  \citet{baumsnow2024} show for the United States
that floor-space supply elasticities (0.5) substantially exceed unit supply
elasticities (0.3), suggesting the quality margin is a first-order phenomenon
that standard count-based measures miss.  Third, it has distributional
implications: if high-walkability neighborhoods are in the upper tail of the
price distribution, restricting their supply concentrates housing services among
high-income households and forces lower-income households into car-dependent,
amenity-poor areas.

The paper proceeds as follows.  Section~\ref{sec:lit} reviews the hedonic price
and effective housing supply literatures.
Section~\ref{sec:context} describes the Israeli housing market context.
Section~\ref{sec:data} describes the urbanism and transaction datasets.
Section~\ref{sec:method} outlines the empirical strategy.
Section~\ref{sec:results} presents the correlation results.
Section~\ref{sec:theory} develops the effective housing supply framework and
applies it to the Israeli case.
Section~\ref{sec:policy} discusses policy implications.
Section~\ref{sec:conclusion} concludes.

%======================================================================
\section{Literature Review}
\label{sec:lit}
%======================================================================

\subsection{The Hedonic Price Framework}
\label{sec:lit_hedonic}

The workhorse model connecting neighborhood characteristics to housing prices is
the hedonic price model (HPM), originating with \citet{rosen1974} and
\citet{harrison1978}.  The HPM treats a dwelling as a bundle of structural
attributes (rooms, age, floor area), locational characteristics (distance to
employment centers, school quality), and neighborhood amenities (parks, crime
rates, environmental quality).  In competitive equilibrium, the market price of
a dwelling equals the sum of implicit prices of its constituent attributes:
\begin{equation}
\label{eq:hedonic}
  P_j = P\bigl(S_j,\, L_j,\, N_j\bigr),
\end{equation}
where $S_j$ denotes structural characteristics, $L_j$ location characteristics,
and $N_j$ neighborhood attributes of dwelling $j$.  Regression analysis
recovers the implicit (``shadow'') price of each attribute \citep{malpezzi2003}.

The theoretical foundation of the HPM connects directly to welfare analysis:
the implicit price of an attribute equals households' marginal willingness to
pay for it.  This makes the HPM the natural tool for measuring how much
households value walkability, street connectivity, and amenity access—which is
precisely the information needed to quality-adjust the housing stock.

Recent advances have extended the HPM in two directions relevant to our study.
First, big-data methods—Google Street View images processed by deep learning
\citep{qiu2022}—capture visual streetscape quality at the pedestrian level,
finding substantial implicit prices for street-level greenness and sky openness.
Second, accessibility indices combining road-network travel times with
employment distributions substantially improve HPM predictive performance:
\citet{madrid2023} find that walk-accessibility indices reduce spatial
autocorrelation in hedonic residuals by 35\% in Madrid.  Both developments
underscore the importance of the network structure of the urban environment as
a determinant of residential value.

\subsection{Street-Network Connectivity and Intersection Density}
\label{sec:lit_network}

Street-network form—the density of intersections, block size, and degree of
grid-like connectivity—has attracted growing attention as a determinant of both
travel behavior and property values.  \citet{barrington2015,barrington2020}
document a century of declining street-network connectivity in US cities,
conceptualizing the phenomenon as path-dependent sprawl: once low-connectivity
infrastructure is built, densification becomes structurally constrained.

Intersection density serves as a proxy for walkability, pedestrian route choice,
and transit efficiency.  Smaller blocks enable more direct pedestrian paths,
support more frequent transit stops, and create more human-scaled environments.
LEED for Neighborhood Development (LEED-ND) uses 140 intersections per square
mile as a design standard for walkable neighborhoods.

The price effects of street connectivity are nuanced.  \citet{matthews2007} find
that street layout affects prices through an interaction of accessibility and
land-use mix: connectivity raises prices in mixed-use contexts but may be
associated with lower prices in purely residential areas where cul-de-sac
privacy is valued \citep{song2003}.  \citet{millardball2022} provides an
economic perspective on street width, estimating that the land value of
residential street rights-of-way totals \$959 billion in 20 large US counties,
and that narrowing over-wide streets could free up land worth over \$100,000
per housing unit in high-cost markets.

\subsection{Local Amenities and Walkability}
\label{sec:lit_amenity}

Amenity density—the concentration of retail, food service, healthcare,
education, and other daily-life destinations within walking distance—is one of
the most robustly valued neighborhood attributes in the hedonic literature.
Walk Score and similar composite indices have been found to carry significant
price premia across diverse metropolitan contexts: \citet{pivo2011},
\citet{li2015}, and a Melbourne study spanning all socioeconomic quintiles all
find positive walkability–price associations, with destination accessibility
driving the effect more than transit access per se.

Environmental quality in the broader sense—greenness, air quality, proximity
to parks—also commands significant premia.  \citet{black1999} demonstrated
using regression discontinuity designs that proximity to higher-performing
schools is substantially capitalized into prices.  \citet{franco2018} estimated
implicit prices for urban greenness in Lisbon using NDVI measures and hedonic
methods.

\subsection{Market Access and Residential Land Values}
\label{sec:lit_ma}

A theoretically grounded synthesis of the accessibility findings comes from the
market access framework derived from general equilibrium trade theory.
\citet{donaldson2016} show that all general equilibrium effects of changes in
transport infrastructure on a location's economy are summarized by changes in
its \textit{market access}:
\begin{equation}
  MA_o = \sum_d \tau_{od}^{-\theta} \cdot Y_d,
\end{equation}
where $\tau_{od}$ is the bilateral trade cost, $\theta$ the trade elasticity,
and $Y_d$ the economic mass of destination $d$.  For residential housing, the
analog is a location's access to employment, consumption amenities, and the
broader urban economy—all discounted by commuting or travel costs determined
by the local street network.

\citet{ahlfeldt2015}, in a celebrated \textit{Econometrica} paper, bring this
framework inside the city using Berlin's division and reunification as exogenous
variation.  They identify substantial elasticities of floor-space prices with
respect to commuter market access, operating through both production
externalities (access to employment density) and residential externalities
(access to residential density).  \citet{heblich2020} obtain analogous results
for London using 19th-century underground railway construction.  The structural
interpretation is clear: a location's street network—its connectivity,
directness, and density—determines its market access, which in turn determines
land values.  Our urbanism metrics are reduced-form proxies for this structural
concept.

\subsection{Housing Quality as a Component of Effective Housing Supply}
\label{sec:lit_effhs}

Most spatial equilibrium models—from the foundational \citet{rosen1979}–\citet{roback1982} framework through \citet{moretti2011} and \citet{hsieh2019}—treat
housing as a scalar: each household consumes a quantity $h$ of ``housing
services'' at a per-unit price $P_i$ in city $i$.  Spatial equilibrium requires
that mobile workers be indifferent across locations:
\begin{equation}
\label{eq:spatial_eq}
  V = \frac{W_i \cdot Z_i}{P_i^\beta} \quad \forall\, i,
\end{equation}
where $W_i$ is the wage, $Z_i$ captures amenities, and $\beta$ is the housing
expenditure share.  The critical implicit assumption is that $P_i$ is a
quality-adjusted price—comparable across locations after abstracting from
quality differences.  In practice, this is a strong assumption.

In labor economics, heterogeneous workers are aggregated into effective labor
supply using relative wages as weights \citep{katz1992changes}:
\begin{equation}
\label{eq:efflabor}
  L^{\text{eff}} = \sum_i \frac{w_i}{\bar{w}} \cdot L_i,
\end{equation}
where $w_i$ is the wage of worker type $i$ and $\bar{w}$ is the average wage.
An exact analogy applies to housing: dwelling units can be aggregated into
``effective housing supply'' using hedonic price weights:
\begin{equation}
\label{eq:effhs}
  H^{\text{eff}} = \sum_j \frac{\hat{p}_j}{\bar{p}} \cdot 1,
\end{equation}
where $\hat{p}_j$ is the hedonic-predicted value of dwelling $j$ (reflecting
all observable quality attributes including neighborhood and street quality)
and $\bar{p}$ is the market average.  This is precisely the housing analog
to computing effective labor in efficiency units \citep{barnett1979}.

\citet{baumsnow2024} provide the most rigorous recent implementation of this
idea for the United States.  Decomposing supply responses across 306 metro
areas into a units margin and a quality margin (floor space per unit,
renovations), they find that floor-space supply elasticities average 0.5 while
unit supply elasticities average only 0.3.  The quality margin—ignored by
count-based measures—accounts for nearly half of the total supply response.
They further find that supply responses exhibit more variation \textit{within}
metro areas than \textit{between} them, highlighting the microgeographic
nature of quality heterogeneity.

\citet{albouy2018} introduce the concept of ``home productivity'' $A_j^Y$—the
efficiency with which land and capital are converted into housing services in
city $j$—into the spatial equilibrium framework.  Better neighborhood quality
(street connectivity, transit, walkability) raises effective home productivity
by increasing the utility yield per dollar of housing expenditure.  A dwelling
in a walkable, transit-accessible neighborhood with good schools provides more
utility per dollar of rent than a physically identical unit in a car-dependent,
amenity-poor suburb: neighborhood quality is thus a form of locational home
productivity.

\citet{hsieh2019}'s influential analysis of housing supply misallocation argues
that supply constraints in high-productivity US cities (New York, San
Francisco) have reduced aggregate growth by preventing labor reallocation.
While their headline estimates have been revised downward by
\citet{greaney2023}, the conceptual mechanism survives: the effective cost of
housing is the \textit{price per unit of service}, not the price per unit of
floor space.  A city that provides the same number of apartments at lower
locational quality has a higher effective cost of living, deterring in-migration
just as higher rents would.

Despite these intellectual ingredients being available, there is no widely
adopted, formalized effective housing supply index.  The gap persists because
constructing such an index requires (a) hedonic estimates of neighborhood
quality attributes, (b) a framework for aggregating them, and (c) embedding the
result in a spatial equilibrium model.  This paper takes a step toward filling
that gap for the Israeli context.

%======================================================================
\section{The Israeli Housing Market: Institutional Context}
\label{sec:context}
%======================================================================

Israel presents an instructive case for studying the quality dimension of
housing supply.  Several features of the Israeli urban system make the
quality–price nexus particularly relevant.

\paragraph{Concentrated urban geography.}  Israel is a small, densely populated
country (9.7 million people in 22,000 km$^2$) with highly concentrated
urbanization.  The Gush Dan metropolitan area (Tel Aviv and surrounding cities)
accounts for roughly 45\% of GDP and contains six of our sample cities.  The
resulting spatial compression means that differences in neighborhood quality
within the urban system are economically large relative to distances.

\paragraph{Rapid price appreciation.}  Real apartment prices roughly doubled
between 2008 and 2023, driven by population growth (roughly 2\% per year),
underbuilding during the 2000s, and low interest rates.  The Israel Land
Authority (ILA) controls roughly 93\% of land, creating structural constraints
on supply that go beyond standard planning regulation.  Housing affordability
has become a central political issue, triggering multiple government plans to
accelerate construction.

\paragraph{Planning and zoning.}  Israeli planning operates under the Planning
and Building Law of 1965 (with major amendments).  Urban areas are governed by
city-level outline plans (toch\-niot m'tah) that specify permitted land uses,
building rights, and density.  In practice, plan approval and building permit
processes are slow, and densification in existing walkable neighborhoods faces
significant opposition.  New construction is often pushed to peripheral areas
with lower land costs but also lower walkability and amenity density.

\paragraph{Urban heterogeneity.}  Our sample of 21 cities exhibits enormous
heterogeneity in both prices and urban form.  Tel Aviv (median price
\textsterling 3.2 million; walkability index 74.6; amenity density 437/km$^2$)
is at one extreme; Be'er Sheva (median price \textsterling 885,000; walkability
index 40.3; amenity density 29.4/km$^2$) at the other.  This variation makes
Israel an ideal laboratory for studying the urbanism–price relationship.

%======================================================================
\section{Data}
\label{sec:data}
%======================================================================

\subsection{Housing Transaction Data}
\label{sec:data_housing}

Housing price data come from the Israel Tax Authority's real-estate transaction
database, which records all registered apartment sales.  We obtained
municipality-level extracts for 20 cities covering transactions from January
2018 through December 2023—the most recent six-year window available.  Each
record includes the deal amount (in NIS), the deal date, the asset type
(apartment, house, commercial, etc.), and the number of rooms.

We restrict the sample to \textit{apartment sales} (Hebrew: \textit{dira}),
filtering on the deal nature description field.  We exclude transactions with
missing or implausible amounts and trim the top and bottom 1\% of prices within
each city to reduce the influence of luxury outliers and data errors.  After
cleaning, our sample comprises approximately 200,000 transactions across 20
cities.  We aggregate to the city level by computing the median deal amount
per city.  (We use the median rather than the mean because apartment prices are
right-skewed; median and mean price correlations with urbanism metrics are
similar.)  Table~\ref{tab:housing} reports summary statistics.

\begin{table}[H]
\centering
\caption{City-Level Housing Price Summary, 2018--2023}
\label{tab:housing}
\small
\begin{tabular}{lrrrr}
\toprule
City & Transactions & Median Price (NIS) & Mean Price (NIS) & Median Price/Room (NIS) \\
\midrule
Tel Aviv--Jaffa       & 21,072 & 3,200,000 & 3,641,425 & 1,046,667 \\
Ra'anana              &  3,312 & 2,580,000 & 2,753,410 &   616,667 \\
Modi'in               &  3,595 & 2,340,000 & 2,449,880 &   565,000 \\
Herzliya              &  3,947 & 2,525,000 & 2,769,953 &   656,250 \\
Jerusalem             & 18,565 & 2,050,000 & 2,299,030 &   600,000 \\
Ramat Gan             & 10,315 & 2,191,000 & 2,339,210 &   666,667 \\
Kfar Sava             &  3,695 & 2,250,000 & 2,375,595 &   559,583 \\
Rishon LeZion         &  9,340 & 1,790,000 & 1,913,487 &   500,000 \\
Petah Tikva           & 13,027 & 1,765,000 & 1,910,413 &   482,500 \\
Netanya               & 11,103 & 1,770,000 & 1,957,783 &   466,667 \\
Bat Yam               &  7,712 & 1,500,000 & 1,620,007 &   530,000 \\
Rehovot               &  6,397 & 1,800,000 & 1,896,672 &   471,429 \\
Holon                 &  8,728 & 1,690,000 & 1,786,862 &   496,667 \\
Bnei Brak             &  6,635 & 1,750,000 & 1,859,413 &   550,000 \\
Ashdod                & 10,711 & 1,600,000 & 1,697,137 &   442,500 \\
Ashkelon              &  7,188 & 1,200,000 & 1,242,278 &   316,667 \\
Beit Shemesh          &  4,891 & 1,485,000 & 1,594,492 &   405,000 \\
Haifa                 & 19,881 & 1,170,000 & 1,317,500 &   350,000 \\
Hadera                &  4,917 & 1,340,000 & 1,375,467 &   347,500 \\
Be'er Sheva           & 14,664 &   885,000 &   972,451 &   266,667 \\
\midrule
Mean (unweighted)     &  9,479 & 1,893,550 & 2,073,079 &   518,920 \\
\bottomrule
\end{tabular}
\end{table}

The price gradient from Tel Aviv (\textsterling 3.2M) to Be'er Sheva
(\textsterling 885K) reflects a 3.6-fold range—enormous for a country the size
of Israel.  This variation is partly compositional (apartment size, age), but
our subsequent analysis shows that it is also strongly associated with
differences in urbanism quality.

\subsection{Urbanism Quality Metrics}
\label{sec:data_urbanism}

Urbanism metrics are computed at the city level using open-source street-network
and point-of-interest (POI) data from OpenStreetMap, retrieved via Overture Maps
(Amazon S3 parquet files, 2024 vintage).  City boundaries are defined using the
Israeli Central Bureau of Statistics 2022 statistical area (SA) geodatabase
(\textit{ezorim statistiim}), projected onto Israeli Transverse Mercator
(EPSG:2039) for area calculations.  We compute six metrics:

\paragraph{Junction density} (intersections/km$^2$).  Nodes in the pedestrian
road graph with degree $\geq 3$, divided by city area.  Higher values indicate
a denser, more grid-like street network.  This metric is the standard empirical
proxy for walkable block structure \citep{barrington2015}.

\paragraph{Street density} (km/km$^2$).  Total length of the walkable road
network divided by city area.  A city with more street kilometres per unit area
provides residents with more route options and shorter detours.

\paragraph{Dead-end ratio}.  The share of nodes with degree 1 (cul-de-sac
termini or stub streets) in total nodes (intersections plus dead ends).  Higher
values indicate a tree-like, disconnected network—the street-network signature
of post-war suburban planning.

\paragraph{Circuity} (dimensionless ratio).  The mean ratio of network distance
to straight-line (Euclidean) distance across all edges.  Values closer to 1.0
indicate straighter, more direct streets; higher values indicate winding,
indirect routes.  Circuity captures a different dimension of connectivity from
intersection density: a highly connected network can still be circuitous if
streets follow topographic contours.

\paragraph{Amenity density} (POIs/km$^2$).  The count of daily-life amenity
POIs—shops, supermarkets, pharmacies, hospitals, schools, restaurants, cafes,
banks, and post offices—divided by city area.  This metric directly captures
the concentration of walkable destinations, the primary driver of walkability
premia found in the hedonic literature \citep{pivo2011}.

\paragraph{Walkability index} (0--100).  A composite score combining the five
network and amenity metrics, min-max normalised and weighted as follows:
30\% junction density, 20\% street density, 15\% dead-end ratio (inverted),
15\% circuity (inverted), 20\% amenity density.  This weighting scheme is
consistent with the evidence that pedestrian connectivity and destination
accessibility are the primary drivers of walkability premia
\citep{bartholomew2011,pivo2011}.

Table~\ref{tab:urbanism} reports summary statistics for the urbanism metrics.
The variation is substantial: junction density ranges from 5.5 (Hadera,
partially a large-area municipality) to 119.1 intersections/km$^2$ (Tel Aviv);
amenity density ranges from 2.9 to 437.4 POIs/km$^2$.

\begin{table}[H]
\centering
\caption{Urbanism Metrics: Summary Statistics ($n=21$ city observations)}
\label{tab:urbanism}
\small
\begin{tabular}{lrrrrrr}
\toprule
Metric & Min & 25th & Median & 75th & Max & Std Dev \\
\midrule
Junction density (int./km$^2$)     &   5.5 &  47.9 &  77.9 & 104.2 & 119.1 & 30.4 \\
Street density (km/km$^2$)         &  12.2 &  17.4 &  24.1 &  31.3 &  33.7 &  7.1 \\
Dead-end ratio (share)             &  0.48 &  0.62 &  0.66 &  0.69 &  0.78 &  0.07 \\
Circuity (ratio)                   &  1.15 &  1.30 &  1.34 &  1.40 &  1.65 &  0.12 \\
Amenity density (POIs/km$^2$)      &   2.9 &  27.0 &  99.4 & 175.0 & 437.4 & 115.2 \\
Walkability index (0--100)         &  19.9 &  41.0 &  59.9 &  70.3 &  74.8 & 15.3 \\
\bottomrule
\end{tabular}
\end{table}

\subsection{Data Linkage}
\label{sec:data_linkage}

The housing price data are organized by municipal code (SEMEL\_YISHUV), while
the urbanism metrics are attached to CBS statistical areas.  We match them by
city name, applying a small set of canonical name normalizations (e.g., ``Tel
Aviv-Jaffa'' to ``Tel Aviv--Jaffa'' in the Overture dataset; ``Herzliya''
spelling variants).  Of the 20 housing-price cities, 21 city-level matches are
obtained (Hadera is matched to two non-overlapping statistical areas and
averaged).  The matched sample of 21 city-area pairs forms the core analytical
dataset.

%======================================================================
\section{Empirical Strategy}
\label{sec:method}
%======================================================================

\subsection{Unit of Analysis and Identification}
\label{sec:method_unit}

Our analysis is at the \textit{city level}: one observation per city, with the
outcome variable being the city's median apartment price and the explanatory
variables being its urbanism metrics.  This choice is dictated by the
availability of the urbanism metrics (computed at the city/SA level from
open-source network data) and the housing price data (available at the
municipality level from administrative records).

We are explicit that city-level correlations are \textit{not} causal estimates.
Unobserved city characteristics—history, wealth, socioeconomic composition,
public investment, regulatory environment—are correlated with both urbanism
metrics and prices.  Tel Aviv's high junction density and high prices both
reflect its position as Israel's economic and cultural center; causality could
run in both directions, or from unobserved third factors.

The cross-city correlations nonetheless serve two purposes.  First, they
establish the empirical pattern: across Israeli cities, urban form and price are
strongly co-determined, consistent with the hedonic literature and with spatial
equilibrium theory.  Second, they provide the raw co-variation needed to
calibrate an effective housing supply index: whether or not walkability
\textit{causes} higher prices, the observed price premia in walkable cities
represent the market's revealed valuation of urban quality—the implicit price
needed for quality adjustment.

\subsection{Correlation Analysis}
\label{sec:method_corr}

For each urbanism metric $x$ and city-level median price $p$, we compute:
\begin{itemize}
  \item \textbf{Pearson correlation} $r$, which measures the linear association.
  \item \textbf{Spearman rank correlation} $\rho$, which is robust to outliers
    and captures monotone but nonlinear associations.
\end{itemize}
Both are tested against the null of zero correlation using $t$-statistics with
$n-2$ degrees of freedom.  With $n=21$ observations, the critical value for
$p<0.05$ is $|r|>0.433$.

We also estimate bivariate OLS regressions of log median price on each metric
to quantify the price elasticities, with and without Hadera (the outlier city
whose urbanism metrics differ markedly due to the large area of one of its
statistical areas).  These regressions are purely descriptive; we do not claim
to identify causal effects.

\subsection{Walkability–Price Gradient}
\label{sec:method_grad}

To illustrate the economic magnitude of the urbanism–price association, we
estimate the price gradient across the walkability distribution using a
nonparametric locally weighted regression (LOWESS) of median price on walkability
index, and report the price differential between the 25th and 75th percentile
cities.

%======================================================================
\section{Results}
\label{sec:results}
%======================================================================

\subsection{Cross-City Correlations}
\label{sec:results_corr}

Table~\ref{tab:corr} reports Pearson and Spearman correlations between each
urbanism metric and city-level median apartment price, along with their
$p$-values and sample sizes.

\begin{table}[H]
\centering
\caption{Correlations between Urbanism Metrics and Median Apartment Price}
\label{tab:corr}
\begin{tabular}{lcccccc}
\toprule
 & \multicolumn{2}{c}{Pearson} & \multicolumn{2}{c}{Spearman} & & \\
\cmidrule(lr){2-3}\cmidrule(lr){4-5}
Metric & $r$ & $p$ & $\rho$ & $p$ & $n$ & Sig. \\
\midrule
Street density (km/km$^2$)     & 0.633 & 0.002 & 0.644 & 0.002 & 21 & *** \\
Amenity density (POIs/km$^2$)  & 0.636 & 0.002 & 0.470 & 0.032 & 21 & **  \\
Junction density (int./km$^2$) & 0.540 & 0.012 & 0.474 & 0.030 & 21 & *   \\
Walkability index (0--100)     & 0.438 & 0.047 & 0.359 & 0.111 & 21 & *   \\
Circuity (ratio)               & 0.299 & 0.187 & 0.232 & 0.312 & 21 & ns  \\
Dead-end ratio (share)         & 0.212 & 0.356 & 0.277 & 0.225 & 21 & ns  \\
\bottomrule
\multicolumn{7}{l}{\footnotesize $^*p<0.05$; $^{**}p<0.01$; $^{***}p<0.001$. All two-tailed tests.}
\end{tabular}
\end{table}

\paragraph{Street density and amenity density.}  The two strongest predictors
are street density and amenity density, with Pearson correlations of 0.633 and
0.636 respectively, both significant at the 1\% level.  The Spearman
correlation for street density (0.644) is essentially identical to the Pearson,
confirming a linear and monotone relationship.  For amenity density, the
Spearman correlation (0.470) is lower than Pearson (0.636), reflecting the
influence of Tel Aviv as a high-leverage observation with both the highest
amenity density and the highest prices.

\paragraph{Junction density.}  Junction density carries a Pearson correlation
of 0.540 and Spearman of 0.474, both significant at the 5\% level.  This is
consistent with the literature on intersection density as a walkability proxy:
cities with more grid-like, connected street networks command higher prices.

\paragraph{Walkability index.}  The composite walkability index is significant
by Pearson ($r=0.438$, $p=0.047$) but not by Spearman ($\rho=0.359$,
$p=0.111$), suggesting the relationship is somewhat sensitive to the
distributional assumptions.  The composite index aggregates five sub-metrics,
some of which (dead-end ratio, circuity) are not individually significant.  The
weaker Spearman performance reflects the non-monotone contribution of these
sub-metrics.

\paragraph{Dead-end ratio and circuity.}  Neither metric is significantly
associated with city-level apartment prices.  This contrasts with street-level
hedonic studies that sometimes find negative effects of dead-end streets.  At
the cross-city level, the dead-end ratio may reflect compositional differences
in urban structure (older cores vs. newer suburbs within the same municipal
boundary) that are masked by city-level aggregation.  Circuity reflects
topographic features (Jerusalem's hills, Haifa's Carmel ridge) that are not
per se negatively valued by residents.

\subsection{Economic Magnitude}
\label{sec:results_magnitude}

To translate correlations into economic magnitudes, we compare cities at the
25th and 75th percentiles of the walkability and amenity distributions.

\textit{Amenity density}.  The 25th-percentile city (around 27 POIs/km$^2$,
e.g., Beit Shemesh) has a median price of approximately \textsterling 1.49M.
The 75th-percentile city (around 175 POIs/km$^2$, e.g., Bnei Brak) has a
median price of approximately \textsterling 1.75M.  Moving from the 25th to the
75th percentile of amenity density is associated with a price premium of
approximately 17\%—roughly \textsterling 260,000 per apartment.

\textit{Street density}.  Moving from the 25th to the 75th percentile of street
density is associated with a price increase of approximately 45\%—from
\textsterling 1.4M to \textsterling 2.0M.  This is a large magnitude reflecting
the fact that street density differences between Israeli cities partly correspond
to the Tel Aviv vs. periphery distinction.

\textit{Walkability index}.  Moving from a walkability index of 41 (Be'er
Sheva, Beit Shemesh) to 70 (Holon, Bat Yam, Ra'anana) is associated with a
price difference of roughly 80--100\%, from \textsterling 885K--\textsterling
1.49M to \textsterling 1.5M--\textsterling 2.6M.  The large magnitude at the
extremes includes differences in city income and demand, but the gradient is
consistent with the walkability capitalization literature.

\subsection{City-Level Patterns}
\label{sec:results_cities}

Figure~\ref{fig:scatter} (see \texttt{data/processed/maps/scatter\_all\_metrics.png})
displays scatter plots of each metric against median price.  Several patterns
are notable:

\begin{itemize}
  \item \textbf{Tel Aviv} is the highest-priced city and also the most walkable
    (junction density 119/km$^2$, amenity density 437/km$^2$, walkability 74.6).
    It is a high-leverage observation but not an outlier—its position is
    consistent with the overall trend.

  \item \textbf{Bnei Brak} is a case study in the limits of physical metrics:
    it has high junction density (111/km$^2$), high amenity density
    (191/km$^2$), and a walkability score of 74.6—very similar to Tel Aviv—yet
    its median price (NIS 1.75M) is barely half of Tel Aviv's.  This reflects
    unobserved socioeconomic factors: Bnei Brak is an ultra-orthodox city with
    a distinct demand structure.

  \item \textbf{Be'er Sheva} is the lowest-priced city and has the lowest
    walkability (40.3) and the second-lowest amenity density (29.4/km$^2$),
    consistent with its role as the regional capital of the Negev desert with
    large low-density residential areas.

  \item \textbf{Haifa} exhibits a paradox: moderate walkability (56.4) but
    among the lowest prices (NIS 1.17M).  Haifa's topography (Carmel ridge)
    makes its street network somewhat circuitous despite reasonable connectivity
    metrics, and the city has experienced long-run population stagnation.
\end{itemize}

These patterns underscore that cross-city urbanism metrics capture an important
dimension of housing value but cannot fully explain price variation—unobserved
city-level fundamentals (income, employment, history, amenity composition)
matter as well.

%======================================================================
\section{Effective Housing Supply: Theory and Israeli Application}
\label{sec:theory}
%======================================================================

\subsection{Formalizing Effective Housing Supply}
\label{sec:theory_formal}

The empirical correlations in Section~\ref{sec:results} establish that the
Israeli housing stock is highly quality-heterogeneous: the same physical
apartment contributes very different amounts of ``housing services'' depending on
whether it is located in Tel Aviv or Be'er Sheva.  We now formalize this
observation using the effective housing supply framework sketched in
Section~\ref{sec:lit_effhs}.

Let $j$ index dwelling units and $P_j$ denote the market transaction price.  In
a competitive market, $P_j$ equals the present discounted value of housing
services that dwelling $j$ provides.  Decompose $P_j$ as:
\begin{equation}
\label{eq:decomp}
  P_j = \underbrace{P^S(S_j)}_{\text{structural}}
      + \underbrace{P^L(L_j)}_{\text{locational}}
      + \underbrace{P^N(N_j)}_{\text{neighborhood}}
      + \varepsilon_j,
\end{equation}
where $P^S$, $P^L$, $P^N$ are the hedonic price functions for structural,
locational, and neighborhood attributes respectively, and $\varepsilon_j$ is a
residual.  The ``structural'' component captures apartment size, age, and
condition—the attributes typically counted in unit-supply statistics.  The
``locational'' and ``neighborhood'' components capture exactly the urbanism
metrics we study: market access, walkability, amenity density.

Define the \textit{locational quality multiplier} for city $i$ as:
\begin{equation}
\label{eq:lqm}
  \lambda_i = \frac{\bar{P}_i}{\bar{P}},
\end{equation}
where $\bar{P}_i$ is the city's median hedonic-predicted price (or, in our
reduced-form context, the median transaction price controlling for structural
characteristics) and $\bar{P}$ is the national median.  Then the effective
housing supply in city $i$ is:
\begin{equation}
\label{eq:ehs_city}
  H_i^{\text{eff}} = \lambda_i \cdot H_i,
\end{equation}
where $H_i$ is the raw count of dwelling units.  An apartment in Tel Aviv
contributes $\lambda_{\text{TA}} > 1$ effective housing units; an apartment in
Be'er Sheva contributes $\lambda_{\text{BS}} < 1$.

\subsection{Calibration for Israel}
\label{sec:theory_calib}

Using our transaction data and the national median price of approximately
NIS 1.6M (the unit-weighted average across our sample), we compute locational
quality multipliers for each city.  Table~\ref{tab:multipliers} reports these
multipliers alongside approximate dwelling stock counts (from CBS 2022).

\begin{table}[H]
\centering
\caption{Locational Quality Multipliers and Effective Housing Supply, Selected Cities}
\label{tab:multipliers}
\small
\begin{tabular}{lrrrrr}
\toprule
City & Median Price & $\lambda_i$ & Dwelling Stock & Raw Units & Eff. Units \\
     & (NIS)       &              & (approx., 000s)  &(index)  & (index) \\
\midrule
Tel Aviv      & 3,200,000 & 2.00 & 240 & 100 & 200 \\
Ra'anana      & 2,580,000 & 1.61 &  55 & 100 & 161 \\
Jerusalem     & 2,050,000 & 1.28 & 330 & 100 & 128 \\
Haifa         & 1,170,000 & 0.73 & 175 & 100 &  73 \\
Be'er Sheva   &   885,000 & 0.55 & 155 & 100 &  55 \\
\midrule
\multicolumn{6}{l}{\footnotesize Note: Dwelling stock figures are approximate.
  $\lambda_i = \bar{P}_i / 1{,}600{,}000$ (national sample median).} \\
\bottomrule
\end{tabular}
\end{table}

The effective supply adjustment is substantial.  Tel Aviv's housing stock
contributes twice as many effective housing units as an equivalently sized
stock in the national median city; Be'er Sheva's stock contributes only 55\% as
many.  If Israel needs, say, 50,000 new effective housing units to close its
affordability gap, it would need \textit{more than} 50,000 physical units if
those units are built in low-quality locations (peripheral cities with low
walkability and amenity density), and could achieve the target with
\textit{fewer} physical units if built in high-quality locations.

\subsection{Connection to Housing Supply Constraints}
\label{sec:theory_supply}

The effective housing supply framework connects naturally to the Hsieh--Moretti
spatial misallocation argument.  In their model, the relevant price for labor
allocation is the per-unit cost of housing services $P_i$ in city $i$—the
effective price, not the raw transaction price.  If housing supply is
constrained in high-quality cities (as it is in Tel Aviv, where land is scarce
and planning restrictions are binding), then the cost of accessing a unit of
effective housing services in Tel Aviv is very high, deterring workers from
locating there and causing spatial misallocation.

Our data directly quantify this constraint.  Tel Aviv's walkability index is 74.6,
amenity density 437/km$^2$—the highest in our sample by a wide margin.  Its
median price is NIS 3.2M, more than three times the price in Be'er Sheva.
The price differential reflects the premium households are willing to pay for
walkable, amenity-rich neighborhoods.  If Tel Aviv could expand its effective
housing supply—either by building more units at current quality levels or by
upzoning to allow denser mixed-use development—the effective price of a housing
service unit would fall, enabling more households to access high-quality urban
environments.

Crucially, however, expanding the \textit{quantity} of units in low-walkability
peripheral cities does not expand effective supply proportionally, because these
units carry lower locational quality multipliers.  The locational quality
margin—neighborhood quality as home productivity—means that building at the
periphery is an imperfect substitute for building in walkable, amenity-rich
areas.

\subsection{The Quality Elasticity}
\label{sec:theory_elasticity}

Extending the \citet{baumsnow2024} framework to include locational quality,
the total supply response to a demand shock $\Delta D$ in city $i$ can be
decomposed as:
\begin{equation}
\label{eq:supply_decomp}
  \frac{\Delta H_i^{\text{eff}}}{\Delta D}
    = \underbrace{\lambda_i \cdot \frac{\partial H_i}{\partial D}}_{\text{units at current quality}}
    + \underbrace{H_i \cdot \frac{\partial \lambda_i}{\partial D}}_{\text{locational quality change}}.
\end{equation}
The first term captures the conventional unit supply elasticity, scaled by the
city's existing quality level.  The second term captures the endogenous
response of neighborhood quality itself to population growth—the agglomeration
of amenities, the upgrading of retail, and the improvement of transit induced
by higher density.  \citet{ahlfeldt2015} identify both production and
residential agglomeration externalities with elasticities of 0.07 and 0.14
respectively; these are precisely the channels through which $\partial \lambda_i /
\partial D > 0$.

For Israel, this decomposition implies that cities with high baseline walkability
(Tel Aviv, Ramat Gan) can expand effective supply both at the extensive margin
(more units) and at the intensive margin (higher walkability from increased
density).  Peripheral cities with low walkability and low amenity density expand
effective supply only at the extensive margin and may see limited locational
quality appreciation even as population grows.

%======================================================================
\section{Policy Implications}
\label{sec:policy}
%======================================================================

Our findings carry several implications for Israeli housing policy.

\paragraph{Quality-adjusted housing targets.}  National housing construction
targets in Israel are typically stated in terms of raw unit counts (e.g., the
government's target of 50,000 new units per year).  Our analysis suggests that
a quality-adjusted target—one that weights new units by their locational
quality multiplier—would better reflect the actual contribution to housing
welfare.  A new apartment in Tel Aviv's Florentine neighborhood contributes
roughly twice as much to effective housing supply as an identical apartment in
a peripheral city.

\paragraph{Location of new supply.}  The strong correlation between walkability
and price implies that there is unmet demand for walkable, amenity-rich
neighborhoods that is not being satisfied by new construction.  If new supply
is systematically directed to low-walkability peripheral areas (as the ILA land
release pattern tends to imply), the effective supply gap may widen even as
raw unit counts increase.  Planning reform should prioritize densification of
existing walkable neighborhoods over greenfield development at the periphery.

\paragraph{Street-network investment as housing policy.}  Our results show that
street density and amenity density are the strongest predictors of apartment
prices.  This implies that investments in street-network improvement—adding
intersections, reducing block sizes, activating ground-floor retail—directly
increase the effective housing supply by raising the locational quality
multiplier of existing units.  Such investments are often cheaper per effective
unit than new construction, and they benefit existing residents as well.

\paragraph{Avoiding the effective supply illusion.}  Israel's central government
has in recent years approved large-scale housing construction in peripheral
cities (Hadera, Ashkelon, Be'er Sheva, Beit Shemesh) on the grounds that this
will reduce affordability pressure in the Tel Aviv metropolitan area.
Our analysis suggests caution: apartments built in cities with walkability
indices of 20--40 provide a fraction of the housing services of Tel Aviv
apartments and are imperfect substitutes for central-city supply.  The
``effective supply illusion''—counting low-quality peripheral units as
equivalent to high-quality central units—may lead policymakers to declare
victory on housing supply while the effective deficit persists.

\paragraph{Implications for the Hsieh--Moretti mechanism.}  Even accepting the
downward-revised estimates of \citet{greaney2023}, the spatial misallocation
mechanism is operative in Israel: high-productivity workers are priced out of
high-walkability, high-productivity cities.  The welfare cost is amplified by
the quality dimension: workers who are forced to live in low-amenity peripheral
cities lose not only proximity to employment but also the daily-life amenities
(walkable retail, schools, parks) that walkable neighborhoods provide.  An
effective housing supply measure that captures this quality dimension would
improve both the measurement and the policy analysis of Israeli housing
misallocation.

%======================================================================
\section{Conclusion}
\label{sec:conclusion}
%======================================================================

This paper has linked quantitative urbanism metrics—street-network density,
junction density, amenity density, and composite walkability—to residential
apartment prices across 21 Israeli cities.  Our main empirical finding is that
street density and amenity density are strongly positively correlated with
city-level median apartment prices (Pearson $r \approx 0.63$--$0.64$,
$p<0.01$), while junction density ($r=0.54$, $p<0.05$) and the composite
walkability index ($r=0.44$, $p<0.05$) are also significant.  Dead-end ratio
and circuity are not significantly associated with prices at the city level.

We then connected these empirical findings to the concept of \textit{effective
housing supply}—the quality-adjusted aggregate stock of housing services,
analogous to efficiency-weighted labor supply in labor economics.  We showed that
the observed price premia in walkable, amenity-rich cities correspond precisely
to the locational quality multipliers needed to aggregate a heterogeneous housing
stock into effective supply units.  Calibrating these multipliers for Israel, we
find that an apartment in Tel Aviv contributes approximately twice as many
effective housing units as a national-median apartment, while an apartment in
Be'er Sheva contributes about half as many.

The implications are substantive.  Israel's housing shortage is not merely a
shortage of physical units; it is a shortage of \textit{effective housing
services}—dwellings in walkable, amenity-rich, well-connected neighborhoods.
Policies that add units in low-walkability peripheral cities provide partial
relief at best, because these units carry low locational quality multipliers.
Effective supply expansion requires either densification of existing walkable
neighborhoods or large-scale investment in the street networks and amenities
that create walkability.

Several limitations of the present analysis should be noted.  First, our
city-level correlations do not establish causality; unobserved city
characteristics confound the urbanism–price relationship.  Future work should
exploit within-city variation in urbanism metrics—ideally using planned
infrastructure improvements as instruments—to identify causal effects.
Second, our walkability index is constructed from open-source network data
without direct measurement of perceived walkability or transit quality.
Third, our effective supply framework uses observed median prices as quality
weights, which embeds all determinants of price—including unobserved
amenities—rather than isolating the contribution of urbanism metrics per se.

Notwithstanding these limitations, the analysis demonstrates that urban form
matters for Israeli housing affordability in ways that conventional unit-count
approaches miss.  As Israel seeks to build its way out of its housing crisis,
the \textit{quality} of where it builds is at least as important as the
\textit{quantity}.

%======================================================================
\section*{Data and Code Availability}
%======================================================================

All data processing and analysis code is available in the project repository.
Urbanism metrics are computed from OpenStreetMap data (open license) retrieved
via Overture Maps.  Housing transaction data are from the Israeli Tax Authority
(publicly available at the open-data portal).  Statistical area boundaries are
from the Israeli Central Bureau of Statistics 2022 geodatabase.

%======================================================================
\bibliography{refs}
\bibliographystyle{aer}
%======================================================================

\end{document}
